\section{模型构建}

\subsection{老年需求空间估计模型}

\subsubsection{人口栅格下推思路}

由于网格尺度老年人口数据通常无法直接获得，本文采用区县老龄化比例、人口栅格和居住 POI 进行下推估计。

\subsubsection{老年需求强度计算}

设网格 $g$ 属于区县 $c$，其老年需求强度为：


\begin{equation}
D_g = \operatorname{Pop}_g \times r_c^{65+} \times W_g
\end{equation}



其中，$\operatorname{Pop}_g$ 是网格人口，$r_c^{65+}$ 是区县 65 岁及以上人口比例，$W_g$ 是居住修正权重。

\subsubsection{居住修正权重设定}

居住修正权重由住宅小区密度、老旧小区密度和建成区强度构成。

说明：

这一节不需要展开过深，重点是说明需求估计不是简单平均分配，而是考虑居住空间差异。


\subsection{养老服务有效供给测度}

\subsubsection{服务设施分类赋权}

将养老服务设施划分为机构养老、社区照护、助餐服务、医养协同和代际共享五类。

建议设定：

\begin{table}[H]
\centering
\caption{养老服务类型与权重含义}
\begin{tabular}{cc}
\toprule
\rowcolor{tablehead}
服务类型 & 权重含义 \\
\midrule
养老机构 & 机构照护能力强 \\
日间照料中心 & 社区照护能力较强 \\
养老服务站 & 居家服务支持 \\
长者食堂 & 助餐支持 \\
医疗与药店 & 医养协同支持 \\
代际共享设施 & 生活与社交支持 \\
\bottomrule
\end{tabular}
\end{table}

\subsubsection{服务供给指数计算}

构建网格服务供给指数：


\begin{equation}
S_g = \sum_{k=1}^{m} w_k X_{gk}
\end{equation}



其中，$X_{gk}$ 是网格 $g$ 中第 $k$ 类服务指标，$w_k$ 是对应权重。

\subsubsection{指标标准化处理}

所有指标统一进行极差标准化或 Z-score 标准化，避免不同量纲影响结果。

说明：

这里不要把 CRITIC-TOPSIS 写太重。供给指数只是为 2SFCA 提供输入，主角不是综合评价。


\subsection{代际共享修正 2SFCA 可达性模型}

\subsubsection{服务点供需比计算}

对每个服务点 $j$，计算其服务范围内的供需比：


\begin{equation}
R_j = \frac{S_j}{\sum_{g \in d_{gj} \leq d_0} D_g f(d_{gj})}
\end{equation}



其中，$S_j$ 是服务点有效供给能力，$D_g$ 是网格需求，$d_{gj}$ 是网格到服务点的出行时间。

\subsubsection{网格服务可达性计算}

对每个网格 $g$，计算其可获得的养老服务水平：


\begin{equation}
A_g = \sum_{j \in d_{gj} \leq d_0} R_j f(d_{gj})
\end{equation}



其中，$A_g$ 越大，说明养老服务可达性越好。

\subsubsection{距离衰减函数设定}

采用高斯距离衰减：


\begin{equation}
f(d)=e^{-(d/\beta)^2}
\end{equation}



说明：

距离越远，服务可获得性越低。
这比简单 15 分钟圈更合理。

\subsubsection{代际共享服务修正}

构建共享服务支撑项：


\begin{equation}
A_g^{*}=A_g+\alpha \operatorname{Shared}_g
\end{equation}



其中，$\operatorname{Shared}_g$ 表示网格代际共享服务水平，$\alpha$ 表示共享服务对养老可达性的缓冲系数。

建议正文设：


\begin{equation}
\alpha = 0.2
\end{equation}



在稳健性中检验：


\begin{equation}
\alpha = 0.1,\ 0.2,\ 0.3
\end{equation}



说明：

这是你们论文的第一个亮点。
它体现了社区养老不是孤立的养老设施问题，而是与社区共享服务体系有关。


\subsection{供需错配指数构建}

\subsubsection{基础错配指数}

构建：


\begin{equation}
\operatorname{Mismatch}_g = Z(D_g) - Z(A_g^{*})
\end{equation}



其中，$Z(D_g)$ 是标准化需求，$Z(A_g^{*})$ 是修正后可达性。

\subsubsection{错配程度分级}

将 $\operatorname{Mismatch}_g$ 按分位数划分为：

\begin{table}[H]
\centering
\caption{错配程度分级标准}
\begin{tabular}{cc}
\toprule
\rowcolor{tablehead}
分位 & 等级 \\
\midrule
前 20\% & 高错配 \\
20\%—80\% & 中错配 \\
后 20\% & 低错配 \\
\bottomrule
\end{tabular}
\end{table}

\subsubsection{供需错配四象限}

根据需求强度和可达性将网格划分为：

\begin{table}[H]
\centering
\caption{供需错配四象限类型}
\begin{tabular}{cc}
\toprule
\rowcolor{tablehead}
类型 & 含义 \\
\midrule
高需求—高供给 & 基本匹配区 \\
高需求—低供给 & 优先补位区 \\
低需求—高供给 & 可能冗余区 \\
低需求—低供给 & 基础薄弱区 \\
\bottomrule
\end{tabular}
\end{table}

说明：

这一节直接服务后面的政策建议。


\subsection{GraphSAGE 空间错配识别模型}

\subsubsection{空间图构建}

将 500\,m 网格视为节点，若两个网格空间相邻，则建立边。

节点特征包括：

\begin{itemize}
\item 老年需求强度；
\item 养老服务可达性；
\item 医养协同水平；
\item 代际共享服务水平；
\item 交通可达性；
\item 社区活力指标。
\end{itemize}

\subsubsection{高错配风险标签构造}

根据 $\operatorname{Mismatch}_g$ 的分位数将网格划分为高、中、低错配风险。

GraphSAGE 的任务是预测网格错配风险等级。

\subsubsection{模型聚合思想}

GraphSAGE 通过聚合邻近网格的特征来更新节点表示：


\begin{equation}
h_v^k = \sigma \left(W^k \cdot \operatorname{CONCAT}\left(h_v^{k-1}, \operatorname{AGG}\{h_u^{k-1}\mid u \in N(v)\}\right)\right)
\end{equation}



说明：

正文中不需要讲太多神经网络原理。
重点写它的作用：

\begin{quote}
GraphSAGE 用于识别考虑邻近社区资源共享后的高错配风险区域。
\end{quote}

\subsubsection{模型输出}

输出：

\begin{itemize}
\item 高错配风险概率；
\item 高风险网格空间分布；
\item 与传统错配指数的对比结果。
\end{itemize}


\subsection{LightGBM-SHAP 成因解释模型}

\subsubsection{被解释变量}

被解释变量为：


\begin{equation}
\operatorname{Mismatch}_g
\end{equation}



或高错配风险等级。

\subsubsection{解释变量}

包括：

\begin{itemize}
\item 老年人口密度；
\item 高龄人口密度；
\item 长者食堂密度；
\item 养老服务站密度；
\item 医疗设施密度；
\item 公共交通可达性；
\item 代际共享服务密度；
\item 道路密度；
\item 社区活力指标。
\end{itemize}

\subsubsection{LightGBM 建模}

使用 LightGBM 拟合非线性关系，识别影响错配的关键变量。

\subsubsection{SHAP 解释}

利用 SHAP 输出变量贡献。

重点分析：

\begin{itemize}
\item 哪些变量推高错配；
\item 哪些变量缓解错配；
\item 不同区域的主要短板是否不同。
\end{itemize}

说明：

这一节主要靠 SHAP 图说话，文字不宜过长。


\subsection{MCLP 服务补位优化模型}

\subsubsection{候选补位点设定}

候选点包括：

\begin{itemize}
\item 社区服务中心；
\item 居委会；
\item 社区卫生服务中心；
\item 现有养老服务站；
\item 长者食堂周边；
\item 高错配网格中心。
\end{itemize}

\subsubsection{最大覆盖目标函数}

目标是在有限新增点位下覆盖最多高需求老人：


\begin{equation}
\max \sum_g D_g y_g
\end{equation}



约束：


\begin{equation}
y_g \leq \sum_{j \in N_g}x_j
\end{equation}




\begin{equation}
\sum_j x_j \leq p
\end{equation}



其中，$x_j=1$ 表示在候选点 $j$ 新建或升级服务点，$y_g=1$ 表示网格 $g$ 被覆盖。

\subsubsection{政策情景设定}

设置三种情景：

\begin{table}[H]
\centering
\caption{MCLP 政策情景设定}
\begin{tabular}{cc}
\toprule
\rowcolor{tablehead}
情景 & 新增 / 升级点位 \\
\midrule
保守情景 & 10 个 \\
中等情景 & 20 个 \\
积极情景 & 30 个 \\
\bottomrule
\end{tabular}
\end{table}

\subsubsection{优化效果评价}

评价指标包括：

\begin{itemize}
\item 高需求老人覆盖率；
\item 平均可达时间；
\item 高错配网格数量；
\item 服务公平性变化。
\end{itemize}

说明：

MCLP 是你们论文的落地模型。
它负责回答“补哪里”。
